
\subsubsection{2.1 Execution Feedback for Iterative Code Repair}

Iterative self-repair using execution feedback --- generate code, execute it, use error output as context, re-generate --- is a well-established approach to inference-time code quality improvement. Chen et al. [2023] demonstrated that execution trace feedback enables LLMs to self-debug, improving MBPP by 12\% and achieving approximately $10\times$ sample efficiency compared to best-of-N sampling. Shinn et al. [2023] showed that verbal reinforcement via execution results achieves 91\% HumanEval pass@1 in the Reflexion framework. Gehring et al. [2024] demonstrated that RL-grounded execution feedback (RLEF) further improves over prompting-only repair, reducing required samples by $10\times$. Arimbur [2026] showed that modern 8B instruction-tuned models, including Llama 3.1 8B, achieve meaningful self-repair improvements (+4.9 to +17.1pp on HumanEval) with execution feedback alone, without fine-tuning, and that most gains occur in repair rounds 1--2.

These works collectively establish that execution feedback produces reliable pass@1 improvements. None, however, compare execution feedback against pylint/mypy static analysis on the same benchmarks under compute-controlled conditions. The present study fills this gap.

\subsubsection{2.2 Static Analysis Feedback for LLM Code Quality}

Pylint, mypy, and related static analysis tools have been integrated into LLM code generation pipelines primarily for quality metrics beyond functional correctness. Blyth et al. [2025] demonstrated that iterative pylint/bandit feedback reduces security issues in LLM-generated code from 40\% to 13\% over 10 iterations on PythonSecurityEval. This result establishes that pylint can provide useful feedback --- but on a benchmark where pylint's Error and Warning categories (security violations) are the dominant failure mode. HumanEval and MBPP failures are dominated by logic and runtime errors rather than security issues, making generalization from Blyth et al.'s setting to functional correctness benchmarks unwarranted without direct measurement.

FeedbackEval \citep{Dai2025} provides the most systematic comparison of feedback types, covering compiler feedback, test feedback, minimal feedback, and LLM-expert feedback across HumanEval, CoderEval, and SWE-bench. FeedbackEval establishes that test feedback outperforms compiler feedback (syntax errors only). However, FeedbackEval excludes semantic static analysis (pylint/mypy-style warnings) from its comparison set --- it uses compiler output rather than pylint-style quality warnings --- and does not control for token budget across conditions. The present study directly addresses this gap by adding pylint/mypy as a treatment condition and holding token budget constant across all conditions.

\subsubsection{2.3 Type-Constrained Decoding}

Mündler et al. [2025] showed that type-constrained decoding --- enforcing type correctness via vocabulary-filtered generation --- reduces compilation errors by over 50\% on HumanEval and MBPP. This approach prevents certain errors at generation time rather than correcting them after generation, making it a prevention strategy rather than a repair strategy. The h-m3 experiment planned for this study --- comparing type-constrained decoding to execution feedback repair --- was not executed due to resource constraints. Results from Mündler et al. [2025] suggest that type errors are a minority of HumanEval failures (consistent with the zero mypy coverage finding reported in Section 5.2), and type-constrained decoding addresses a different failure mode than execution feedback repair.

\subsubsection{2.4 Positioning}

This study differs from prior work in three ways: first, it provides a head-to-head comparison of pylint/mypy versus execution feedback on the same benchmarks; second, it enforces iso-compute control (B=1000 output tokens per problem across all conditions); and third, it includes mechanism analysis decomposing pylint coverage by flag category rather than treating total coverage as a proxy for feedback informativeness.

